\newpage
\section{The Build System (CMake)}

FP-Tools uses CMake as its build system. This section covers the reasoning
behind that choice, the required dependencies, the available build options, and
step-by-step build instructions. A detailed walkthrough of the 
\texttt{CMakeLists.txt} file and dependency files themselves are also provided 
for individuals who want to understand or modify the build system.

\subsection{Why CMake?}

CMake is a cross-platform build system generator that's pretty widely used in
C++ projects. Rather than needing to specify how to build the code directly, 
CMake produces native build files for whatever system you are on from a single,
platform-agnostic, description of the build. This makes it very straightforward 
for users to install FP-Tools, for example, on either some personal laptop or
on an HPC machine without maintaining a separate build configuration for each.
\vspace{1em}

Beyond being extremely portable, CMake has great built-in support for finding
and linking external dependencies such as FFTW, OpenMP, and MPI (what we need
for this project!). 


\subsection{Requirements and Dependencies}

FP-Tools is supported on Linux, macOS, and HPC systems running Linux. Windows
is not currently supported. The following table summarizes all of the 
required and optional dependencies.
\vspace{1em}

\begin{table}[h]
\centering
\begin{tabular}{llll}
\toprule
Dependency & Version & Required & Notes \\
\midrule
C++ compiler    & C++17 or later  & Yes & GCC $\geq$ 8, Clang $\geq$ 5, 
                                          Intel $\geq$ 19 \\
CMake           & $\geq$ 3.24     & Yes & \\
FFTW3           & $\geq$ 3.3      & Yes & Required for all analyses \\
OpenMP          & ---             & No  & Enables shared-memory parallelism \\
MPI             & ---             & No  & Enables distributed parallelism \\
\bottomrule
\end{tabular}
\label{tab:dependencies}
\end{table}

\subsubsection{Compiler}

FP-Tools requires a C++17-compatible compiler. The following compilers have 
been tested and should work:
\begin{itemize}
    \item GCC $\geq$ 8 -- recommended for Linux and HPC systems
    \item Clang $\geq$ 5  -- recommended for macOS systems
    \item Intel C++ Compiler $\geq$ 19
\end{itemize}

\begin{note}
    On macOS, the system \texttt{clang} that's provided by XCode's Command Line
    Tools should be sufficient. On HPC systems, available compilers and their
    versions can vary, so I'd recommend either checking with your system
    administrator or load an appropriate module, for example:
    \begin{verbatim}
        module load gcc/11.2.0
    \end{verbatim}
\end{note}

If you need more detailed installation instructions or encounter any issues that
may require troubleshooting, see the installation documentation for each 
respective compiler: 
\href{https://gcc.gnu.org/wiki/InstallingGCC}{GCC}, 
\href{https://clang.llvm.org/get_started.html}{Clang},
\href{https://www.intel.com/content/www/us/en/developer/tools/oneapi/dpc-compiler-download.html}{Intel}.

\subsubsection{CMake}

CMake version 3.24.0 or later is required. This version included improvements for \texttt{FetchContent} and dependency handling that FP-Tools relies on. You 
can install CMake via your package manager:

\begin{bashcode}
# Linux (apt)
sudo apt install cmake

# macOS (Homebrew)
brew install cmake
\end{bashcode}

On HPC machines where you do not have root access, CMake can be installed to a 
local directory or loaded as a module:

\begin{bashcode}
module load cmake/3.24.0
\end{bashcode}

If you need more detailed installation instructions or encounter any issues that
may require troubleshooting, see the CMake installation documentation 
\href{https://cmake.org/download/}{here}.

\subsubsection{FFTW3}

FFTW3 (Fastest Fourier Transform in the West) is required for all analyses in
FP-Tools. It provides the FFT routines used for correlation function
calculations and, in future versions, for MSD computation. FFTW3 can be
installed as follows:

\begin{bashcode}
# Linux (apt)
sudo apt install libfftw3-dev

# macOS (Homebrew)
brew install fftw
\end{bashcode}

On HPC systems, FFTW3 should be available as a module:

\begin{bashcode}
module load fftw/3.3.10
\end{bashcode}

If not, you can install FFTW3 without administrator permissions by following the
installation instructions \href{https://www.fftw.org/fftw2_doc/fftw_6.html}{here}
and installing the package to a local directory. If FFTW3 is installed in a
non-standard location, you can point CMake to it using:

\begin{bashcode}
cmake -B build -DFFTW3_ROOT=/path/to/fftw3
\end{bashcode}


\subsubsection{OpenMP (Optional)}

OpenMP enables shared-memory parallelism within a single node. It's supported by
all major C++ compilers and shouldn't require any separate installation. You can
enable it at compile time via a CMake flag (see Section(REF TO CMAKE OPTIONS)). 
On macOS, the Apple Clang compiler doesn't include OpenMP support by default. If
you want OpenMP on macOS, install LLVM via Homebrew:

\begin{bashcode}
brew install llvm libomp
\end{bashcode}

\subsubsection{MPI (Optional)}

MPI enables distributed-memory parallelism across multiple nodes and is required 
for large-scale HPC runs. Any MPI implementation is supported, including OpenMPI 
and MPICH. You can install these using your package manager:

\begin{bashcode}
# Linux (apt)
sudo apt install libopenmpi-dev

# Linux (dnf)  
sudo dnf install openmpi-devel

# macOS (Homebrew)
brew install open-mpi
\end{bashcode}

On HPC systems, MPI is almost always available as a module and should be loaded 
before building the code:

\begin{verbatim}
module load openmpi/4.1.2
\end{verbatim}

\subsection{A Comprehensive List of Compilation Flags}

\subsection{Detailed Build Instructions}

\subsection{Troubleshooting Common Errors}

\subsection{The Build System in Detail}

\subsubsection{CMakeLists.txt}
\subsubsection{FPToolsFunctions.cmake}
\subsubsection{FFTW.cmake}
\subsubsection{OpenMP.cmake}
\subsubsection{MPI.cmake}
\subsubsection{tests/CMakeLists.txt}
